% Generated draft: requires mathematical and source review.
\documentclass[11pt]{article}
\usepackage[margin=25mm]{geometry}
\usepackage{fontspec}
\setmainfont{Latin Modern Roman}
\usepackage{amsmath,amssymb,mathtools,unicode-math}
\setmathfont{Latin Modern Math}
\usepackage{enumitem,needspace,xurl,hyperref}
\setlength{\parindent}{0pt}
\setlength{\parskip}{5pt}
\setlength{\emergencystretch}{3em}
\newcommand{\N}{\mathbb{N}}
\newcommand{\Z}{\mathbb{Z}}
\newcommand{\Q}{\mathbb{Q}}
\newcommand{\R}{\mathbb{R}}
\newcommand{\C}{\mathbb{C}}
\newcommand{\F}{\mathbb{F}}
\newcommand{\A}{\mathbb{A}}
\newcommand{\PP}{\mathbb{P}}
\newcommand{\E}{\mathbb{E}}
\newcommand{\eps}{\varepsilon}
\newcommand{\dd}{\,\mathrm{d}}
\begin{document}
\section*{UnsolvedMath ID: 30003578}
\textbf{Source problem code:} catalogue entry
\subsection{Definitions and mathematical statement}
A Gaussian prime is a prime element of the ring \(\Z[i]\). For \(z=a+bi\), its norm is \(\operatorname{Nm}(z)=a^2+b^2\), and its argument is naturally defined modulo \(2\pi\). Multiplication by a unit changes the argument by a multiple of \(\pi/2\). Thus counting prime elements, representatives modulo units, or prime ideals requires a specified convention.

Let \(\Theta_X\) denote the finite multiset of Gaussian-prime angles intended by the source, with norm at most \(X\), and put \(N=N(X)=\#\Theta_X\). The source does not specify the convention defining this multiset, including its treatment of associates, conjugates, multiplicities, and primes on the coordinate axes. Its original formulation uses \(p\leq X\) without defining \(p\); the cleaned statement instead specifies a norm bound.

The following framework makes explicit the additional choices needed to interpret the variance. Given an angular circle \(T_L=\R/L\Z\) of period \(L>0\), a multiset \(\Theta_X\subset T_L\), and a measurable arc \(I_K\subset T_L\), define
\[
N_{K,X}(\theta)=\sum_{\alpha\in\Theta_X}\mathbf{1}_{\theta+I_K}(\alpha),\qquad \theta\in T_L,
\]
where the sum respects multiplicity. If the averaging measure is uniform angular measure, define
\[
\mu_{K,X}=\frac1L\int_{T_L}N_{K,X}(\theta)\,\dd\theta,
\qquad
\operatorname{Var}_{\theta}(N_{K,X})
=\frac1L\int_{T_L}\bigl(N_{K,X}(\theta)-\mu_{K,X}\bigr)^2\,\dd\theta.
\]
Writing \(|I_K|\) for arc length gives \(\mu_{K,X}=N|I_K|/L\). In particular, the mean equals \(N/K\) if \(|I_K|=L/K\). This exact arc normalization and the uniform averaging convention are possible completions of the source, rather than specifications supplied by it; merely saying that an arc has scale \(1/K\) does not determine them.

With the intended conventions and admissible dependence \(K=K(X)\) still to be specified, the supplied Gaussian-prime variance conjecture asks whether
\[
\frac{\operatorname{Var}_{\theta}(N_{K,X})}{N(X)/K}
\sim
\min\!\left\{1,\frac{2\log K}{\log N(X)}\right\}
\qquad (X\to\infty).
\]
Here \(A(X)\sim B(X)\) means \(A(X)/B(X)\to1\) where the ratio is defined. The source gives no domain or growth range for \(K\), no uniformity assertion, no exact arc length or endpoint convention, and no averaging measure for the variance. Consequently, it does not determine a fully quantified conjecture. The displayed relation is retained as the supplied target, subject to these unresolved qualifications.
\subsection{Short English statement}
The supplied conjecture predicts the variance of the number of Gaussian-prime angles in a short moving arc, normalized by \(N/K\). A precise statement requires the prime-counting convention, angular period, exact arc normalization, variance measure, and allowed growth of \(K\). The supplied open label and literature assessment are dated curation, not independently verified conclusions.
\subsection{Sources}
\begin{itemize}
\item Supplied local UnsolvedMath registry record, permanent ID \(30003578\), title \emph{Variance of Gaussian Prime Angles in Short Arcs}, exact source problem code \texttt{catalogue entry}, in the Oberwolfach Reports Open Problems collection. No individual proposer is supplied; the proposed year is recorded as 2017.
\item Supplied report citation: \emph{Automorphic Forms and Arithmetic (2018)}, Oberwolfach Reports, pp. 2475--2537, DOI \texttt{10.4171/owr/2017/40}. \url{https://doi.org/10.4171/owr/2017/40}. The background attributes the problem to workshop open-problem discussions documented in this report. The supplied title year, proposed year, and DOI are retained as recorded.
\item Rudnick--Waxman, supplied literature-source label and link: \url{https://doi.org/10.1007/s11856-019-1867-5}. No article title or more detailed bibliographic attribution is supplied.
\item Huang--Liu--Rudnick, supplied literature-source label and link: \url{https://doi.org/10.4064/aa190331-23-7}. No article title or more detailed bibliographic attribution is supplied.
\item OWR report, also listed among the supplied literature sources: \url{https://doi.org/10.4171/owr/2017/40}.
\item Literature assessment checked at \texttt{2026-08-21}: the supplied curation labels the conjecture open and states that unconditional work establishes only partial almost-all-sector results and related function-field analogues. This is supplied dated curation derived from the OWR record, not an independently verified assessment. The supplied background explicitly requests independent verification of time-sensitive status and cited claims.
\end{itemize}
\end{document}
